\section{Optimal-transport rescue inside the filter} \label{sec:ot-rescue}

The Liu--West shrinkage of Section~\ref{sec:liu-west} keeps the inner
state cloud well-behaved most of the time. Occasionally --- under a
particularly informative observation, or when the parameter cloud
visits a low-likelihood region during HMC exploration --- the inner
filter's effective sample size collapses despite the smoothed
resampling. The framework includes a low-rank Sinkhorn optimal-transport
``rescue'' kernel that fires when this happens, blending a transport
plan into the systematic resample by a sigmoid weight on the ESS
ratio. This section explains what it does and why it has the form it
does.

\subsection{When SIR plus Liu--West is not enough}

After the Liu--West shrinkage step of Section~\ref{sec:liu-west}, the
inner state cloud's effective sample size $N_{\rm eff}$ is monitored
once per bin. If it falls below an OT-rescue threshold (default
$N_{\rm eff} / N < 0.05$, controlled by \texttt{ot\_ess\_frac}), the
cloud is judged degenerate enough that systematic resampling will
produce too many duplicates and the filter's ability to track the
state will be lost. The rescue is a soft, sigmoid-blended OT
resampling rather than a hard branch.

The blending weight is
\begin{equation}
w_{\rm OT} \;=\; w_{\rm OT}^{\max} \;\cdot\;
   \sigmoid\!\left(\frac{\text{ess\_thr} - N_{\rm eff}/N}
                          {\text{ess\_temp}}\right),
\label{eq:ot-blend}
\end{equation}
with $w_{\rm OT}^{\max}$ a small cap (default $0.01$, so even at full
saturation only 1\% of the resample comes from OT),
$\text{ess\_thr} = \texttt{ot\_ess\_frac}$, and
$\text{ess\_temp} = \texttt{ot\_temperature}$. The output of the
resampling step at each bin is then
\begin{equation}
X^{\rm out} \;=\; (1 - w_{\rm OT})\, X^{\rm sys+LW}
                \;+\; w_{\rm OT}\, X^{\rm OT},
\end{equation}
i.e.\ a convex combination of the standard systematic+Liu--West output
and the OT-resampled output. Because both inputs are
$\Pcal$-preserving and the weight is in $[0, w_{\rm OT}^{\max}]$, the
combination is also $\Pcal$-preserving.

The continuous-blend formulation matters for the same reason
ESS-scaled bandwidth matters in Section~\ref{sec:liu-west}: the outer
SMC$^2$ HMC kernel (Section~\ref{sec:smc2}) needs the inner-PF
log-likelihood to be a smooth function of $\thetadyn$. A discontinuous
ESS-threshold branch would create gradient discontinuities at every
particle that crosses the threshold, sabotaging the HMC step-size
adaptation.

\subsection{The OT step itself --- low-rank Sinkhorn} \label{sec:sinkhorn}

The OT-rescue kernel takes the current particle cloud
$\{X^{(i)}\}_{i=1}^{N}$ with weights $\{W^{(i)}\}$ and produces a new
cloud whose empirical distribution approximates the same distribution
as systematic resampling, but via a transport plan rather than a
multinomial draw. The transport plan's smoothness with respect to
$\thetadyn$ is what makes the gradient go through.

\paragraph{Step 1 --- low-rank kernel factorisation.}
Pick $r$ anchor particles by index (\texttt{anchor\_idx}, sampled
uniformly without replacement) and form the Gaussian kernel between
all $N$ particles and these $r$ anchors:
\begin{equation}
K_{NR}^{(i, j)} \;=\;
   \exp\!\left(-\frac{\|X^{(i)} - X^{({\rm anchor}\,j)}\|^2}
                       {2\,\epsilon}\right),
\qquad i = 1, \ldots, N;\, j = 1, \ldots, r,
\end{equation}
column-normalised for stability. This is the Nystr\"om approximation
of the full $N \times N$ Gaussian kernel by $K_{NR} K_{NR}^\top$. Cost:
$O(N r)$ in time and memory. Implementation:
\path{smc2fc/filtering/transport_kernel.py:compute_kernel_factor}.

The low-rank form is essential because the full $N \times N$ kernel,
combined with Sinkhorn's $O(N^2)$ matrix-vector products, would
exhaust GPU memory at the rank-$\sim$400 cloud sizes the framework
uses. Empirical default $r = 5$ (\texttt{ot\_rank}); tested up to
$r = 50$ on tighter problems.

\paragraph{Step 2 --- Sinkhorn iterations in scaling form.}
Solve the entropy-regularised optimal transport problem with target
marginals $a$ (current importance weights) and $b$ (uniform weights):
find scalings $u, v \in \R^N_{>0}$ such that
\begin{equation}
\mathrm{diag}(u)\, K_{\rm approx}\, \mathrm{diag}(v)
\quad\text{has marginals}\quad (a,\, b),
\qquad K_{\rm approx} := K_{NR} K_{NR}^\top.
\end{equation}
Sinkhorn alternates the row-marginal and column-marginal projections
\begin{equation}
u^{(\ell+1)} = a / (K_{\rm approx} v^{(\ell)}),
\qquad
v^{(\ell+1)} = b / (K_{\rm approx}^\top u^{(\ell+1)}),
\end{equation}
each iteration costing $O(N r)$ via $K_{\rm approx} v = K_{NR}
(K_{NR}^\top v)$. Default $\ell_{\max} = 2$ (\texttt{ot\_n\_iter}).
Implementation: \path{smc2fc/filtering/sinkhorn.py:sinkhorn_scalings}.

A subtle but important detail: the iteration is implemented via
\texttt{jax.lax.fori\_loop}, not \texttt{lax.scan}. The reason
(documented at the top of \path{sinkhorn.py}) is that \texttt{lax.scan}
inside a particle-filter scan branch causes XLA to unroll all
inner iterations, leading to compilation-time blow-up. \texttt{fori\_loop}
keeps the inner iteration as a real on-device loop.

\paragraph{Step 3 --- barycentric projection.}
The transport plan
$\Gamma = \mathrm{diag}(u)\, K_{\rm approx}\, \mathrm{diag}(v)$
implicitly defines a coupling between source and target particles. The
new particle positions are computed as the barycentric projection
\begin{equation}
X_i^{\rm OT} \;=\;
   \frac{u_i \, [K_{\rm approx} (v \odot X)]_i}
        {u_i \, [K_{\rm approx}\, v]_i}.
\label{eq:barycentric}
\end{equation}
This is a convex combination of source positions, so it preserves any
$[0, 1]^d$ box constraints automatically (no logit / sigmoid needed
for state-bounded models). Implementation:
\path{smc2fc/filtering/project.py:barycentric_projection}.

\subsection{The full OT-rescue API}

The three steps above are combined in
\path{smc2fc/filtering/resample.py:ot_resample_lr}, which is the only
function the inner PF actually calls. Its signature:
\begin{quote}\small
\verb|ot_resample_lr(particles, log_weights, rng_key, stochastic_indices,|\\
\verb|               epsilon=0.5, n_iter=10, rank=50) -> (N, n_states)|
\end{quote}
The \texttt{stochastic\_indices} argument tells the OT kernel which
state coordinates are subject to noise (and thus need rescuing);
deterministic state components pass through unchanged. The function is
fully differentiable through \texttt{jax.grad} --- this is the key
practical feature, distinguishing it from systematic resampling which
has zero gradient with respect to particle positions.

\subsection{Implementation pointer}

\begin{itemize}
  \item \path{smc2fc/filtering/transport_kernel.py} ($\sim$75 lines):
        Nystr\"om kernel factorisation, factor matvec.
  \item \path{smc2fc/filtering/sinkhorn.py} ($\sim$90 lines): low-rank
        Sinkhorn iterations via \texttt{fori\_loop}.
  \item \path{smc2fc/filtering/project.py} ($\sim$45 lines): barycentric
        projection.
  \item \path{smc2fc/filtering/resample.py} ($\sim$100 lines): the
        public \texttt{ot\_resample\_lr} API.
  \item \path{smc2fc/filtering/gk_dpf_v3_lite.py:50} imports
        \texttt{ot\_resample\_lr}; the sigmoid-blend interpolation
        between systematic+LW and OT outputs is in the inner-PF body.
\end{itemize}

\subsection{Practical takeaway}

The OT rescue is a \emph{safety net}, not the main rejuvenation
mechanism. With the default \texttt{ot\_max\_weight = 0.01}, the OT
contribution to the final particle cloud is tiny in normal operation
and only becomes appreciable when the cloud is genuinely degenerate.
Tuning advice: leave the OT-rescue defaults alone unless you see
inner-PF log-likelihood NaNs or sustained $N_{\rm eff} < 5\%$ across
many bins. If you do, the lever is \texttt{ot\_max\_weight} (raise to
0.05--0.10) and \texttt{ot\_n\_iter} (raise to 5--10).
